\documentclass[12pt,a4paper]{article}
\usepackage[UTF8]{ctex}
\usepackage{amsmath,amssymb,amsthm}
\usepackage{geometry}

\geometry{margin=2.5cm}

\title{小位移线性化：将$\ddot{x}$, $\dot{x}$, $x$替换为$\Delta \ddot{x}$, $\Delta \dot{x}$, $\Delta x$}
\author{第11章补充说明}
\date{}

\begin{document}

\maketitle

\section{问题提出}

在阻抗控制中，我们经常看到这样的表述：

"We could also replace $\ddot{x}$, $\dot{x}$, and $x$ with small displacements $\Delta \ddot{x}$, $\Delta \dot{x}$, and $\Delta x$ from reference values in a controlled motion of the robot."

\textbf{问题}：
\begin{enumerate}
\item 什么是参考值（reference values）？
\item 如何定义小位移$\Delta \ddot{x}$, $\Delta \dot{x}$, $\Delta x$？
\item 如何进行替换？
\item 为什么这样做？
\end{enumerate}

\section{基本概念}

\subsection{参考轨迹}

\textbf{参考值}：在受控运动中，机器人沿着期望的参考轨迹运动。

\textbf{参考轨迹}：
\begin{itemize}
\item $x_r(t)$：参考位置（期望位置）
\item $\dot{x}_r(t)$：参考速度（期望速度）
\item $\ddot{x}_r(t)$：参考加速度（期望加速度）
\end{itemize}

\textbf{物理意义}：
\begin{itemize}
\item 这是机器人应该遵循的理想轨迹
\item 例如：沿着直线移动，或沿着圆弧移动
\end{itemize}

\subsection{实际状态}

\textbf{实际状态}：
\begin{itemize}
\item $x(t)$：实际位置
\item $\dot{x}(t)$：实际速度
\item $\ddot{x}(t)$：实际加速度
\end{itemize}

\textbf{物理意义}：
\begin{itemize}
\item 这是机器人实际的状态
\item 由于干扰、误差等原因，可能与参考轨迹不同
\end{itemize}

\subsection{小位移（扰动）}

\textbf{定义}：小位移是实际状态与参考状态的差值。

\begin{align}
\Delta x(t) &= x(t) - x_r(t) \quad \text{（位置扰动）} \\
\Delta \dot{x}(t) &= \dot{x}(t) - \dot{x}_r(t) \quad \text{（速度扰动）} \\
\Delta \ddot{x}(t) &= \ddot{x}(t) - \ddot{x}_r(t) \quad \text{（加速度扰动）}
\end{align}

\textbf{物理意义}：
\begin{itemize}
\item $\Delta x$：位置误差（偏离参考轨迹的距离）
\item $\Delta \dot{x}$：速度误差（偏离参考速度的大小）
\item $\Delta \ddot{x}$：加速度误差（偏离参考加速度的大小）
\end{itemize}

\section{替换过程}

\subsection{原始方程}

\textbf{阻抗控制方程}：
\begin{equation}
M\ddot{x} + B\dot{x} + Kx = f_{\text{ext}} \tag{11.64}
\end{equation}

其中：
\begin{itemize}
\item $x, \dot{x}, \ddot{x}$：实际状态
\item $M, B, K$：虚拟质量、阻尼、刚度矩阵
\item $f_{\text{ext}}$：外力
\end{itemize}

\subsection{替换步骤}

\textbf{步骤1：定义参考状态}

假设机器人沿着参考轨迹$x_r(t)$运动，参考状态满足：
\begin{equation}
M\ddot{x}_r + B\dot{x}_r + Kx_r = f_{\text{ext},r}
\end{equation}

其中$f_{\text{ext},r}$是参考外力（如果没有外力，$f_{\text{ext},r} = 0$）。

\textbf{步骤2：定义小位移}

\begin{align}
x &= x_r + \Delta x \\
\dot{x} &= \dot{x}_r + \Delta \dot{x} \\
\ddot{x} &= \ddot{x}_r + \Delta \ddot{x}
\end{align}

\textbf{步骤3：代入原始方程}

将替换代入方程(11.64)：
\begin{align}
M(\ddot{x}_r + \Delta \ddot{x}) + B(\dot{x}_r + \Delta \dot{x}) + K(x_r + \Delta x) &= f_{\text{ext}} \\
M\ddot{x}_r + M\Delta \ddot{x} + B\dot{x}_r + B\Delta \dot{x} + Kx_r + K\Delta x &= f_{\text{ext}}
\end{align}

\textbf{步骤4：整理}

将参考项和扰动项分开：
\begin{align}
(M\ddot{x}_r + B\dot{x}_r + Kx_r) + (M\Delta \ddot{x} + B\Delta \dot{x} + K\Delta x) &= f_{\text{ext}} \\
f_{\text{ext},r} + (M\Delta \ddot{x} + B\Delta \dot{x} + K\Delta x) &= f_{\text{ext}}
\end{align}

\textbf{步骤5：得到扰动方程}

定义外力扰动：$\Delta f_{\text{ext}} = f_{\text{ext}} - f_{\text{ext},r}$

得到扰动方程：
\begin{equation}
M\Delta \ddot{x} + B\Delta \dot{x} + K\Delta x = \Delta f_{\text{ext}} \tag{11.64-perturbation}
\end{equation}

\subsection{最终结果}

\textbf{扰动方程}：
\begin{equation}
M\Delta \ddot{x} + B\Delta \dot{x} + K\Delta x = \Delta f_{\text{ext}}
\end{equation}

\textbf{与原始方程对比}：

\begin{center}
\begin{tabular}{|l|l|}
\hline
\textbf{原始方程} & \textbf{扰动方程} \\
\hline
$M\ddot{x} + B\dot{x} + Kx = f_{\text{ext}}$ & $M\Delta \ddot{x} + B\Delta \dot{x} + K\Delta x = \Delta f_{\text{ext}}$ \\
\hline
$x, \dot{x}, \ddot{x}$：实际状态 & $\Delta x, \Delta \dot{x}, \Delta \ddot{x}$：扰动 \\
\hline
$f_{\text{ext}}$：总外力 & $\Delta f_{\text{ext}}$：外力扰动 \\
\hline
\end{tabular}
\end{center}

\section{为什么这样做？}

\subsection{线性化}

\textbf{目的}：将非线性系统线性化，简化分析。

\textbf{假设}：
\begin{itemize}
\item 扰动$\Delta x$, $\Delta \dot{x}$, $\Delta \ddot{x}$很小
\item 系统在参考轨迹附近运动
\item 可以忽略高阶项
\end{itemize}

\textbf{优点}：
\begin{itemize}
\item 扰动方程是线性的（即使原始系统可能是非线性的）
\item 可以使用线性系统理论分析
\item 更容易设计和分析控制器
\end{itemize}

\subsection{小扰动分析}

\textbf{应用场景}：
\begin{itemize}
\item 分析系统对扰动的响应
\item 设计控制器使扰动最小化
\item 评估系统的稳定性
\end{itemize}

\textbf{物理意义}：
\begin{itemize}
\item 如果扰动方程稳定，系统会回到参考轨迹
\item 如果扰动很小，系统性能良好
\end{itemize}

\section{具体例子}

\subsection{例子1：跟踪直线轨迹}

\textbf{参考轨迹}：
\begin{itemize}
\item $x_r(t) = v_0 t$（匀速直线运动）
\item $\dot{x}_r(t) = v_0$（常数速度）
\item $\ddot{x}_r(t) = 0$（零加速度）
\end{itemize}

\textbf{实际状态}：
\begin{itemize}
\item $x(t) = v_0 t + \Delta x(t)$
\item $\dot{x}(t) = v_0 + \Delta \dot{x}(t)$
\item $\ddot{x}(t) = \Delta \ddot{x}(t)$
\end{itemize}

\textbf{扰动方程}：
\begin{equation}
M\Delta \ddot{x} + B\Delta \dot{x} + K\Delta x = \Delta f_{\text{ext}}
\end{equation}

\textbf{物理意义}：
\begin{itemize}
\item 如果$\Delta f_{\text{ext}} = 0$（无外力扰动），且初始扰动很小
\item 系统会逐渐回到参考轨迹：$\Delta x \to 0$
\end{itemize}

\subsection{例子2：跟踪圆弧轨迹}

\textbf{参考轨迹}：
\begin{itemize}
\item $x_r(t) = R\cos(\omega t)$（圆弧运动）
\item $\dot{x}_r(t) = -R\omega\sin(\omega t)$
\item $\ddot{x}_r(t) = -R\omega^2\cos(\omega t)$
\end{itemize}

\textbf{实际状态}：
\begin{itemize}
\item $x(t) = R\cos(\omega t) + \Delta x(t)$
\item $\dot{x}(t) = -R\omega\sin(\omega t) + \Delta \dot{x}(t)$
\item $\ddot{x}(t) = -R\omega^2\cos(\omega t) + \Delta \ddot{x}(t)$
\end{itemize}

\textbf{扰动方程}：
\begin{equation}
M\Delta \ddot{x} + B\Delta \dot{x} + K\Delta x = \Delta f_{\text{ext}}
\end{equation}

\textbf{注意}：即使参考轨迹是时变的，扰动方程的形式相同。

\section{在控制律中的应用}

\subsection{原始控制律}

\textbf{阻抗控制律}（方程11.65）：
\begin{equation}
\tau = J^T(\theta)\left(\tilde{\Lambda}(\theta)\ddot{x} + \tilde{\eta}(\theta, \dot{x}) - (M\ddot{x} + B\dot{x} + Kx)\right)
\end{equation}

\subsection{使用小位移替换}

\textbf{替换}：
\begin{align}
x &= x_r + \Delta x \\
\dot{x} &= \dot{x}_r + \Delta \dot{x} \\
\ddot{x} &= \ddot{x}_r + \Delta \ddot{x}
\end{align}

\textbf{控制律变为}：
\begin{align}
\tau &= J^T(\theta)\Big[\tilde{\Lambda}(\theta)(\ddot{x}_r + \Delta \ddot{x}) + \tilde{\eta}(\theta, \dot{x}_r + \Delta \dot{x}) \\
&\quad - (M(\ddot{x}_r + \Delta \ddot{x}) + B(\dot{x}_r + \Delta \dot{x}) + K(x_r + \Delta x))\Big]
\end{align}

\textbf{展开}：
\begin{align}
\tau &= J^T(\theta)\Big[\tilde{\Lambda}(\theta)\ddot{x}_r + \tilde{\Lambda}(\theta)\Delta \ddot{x} + \tilde{\eta}(\theta, \dot{x}_r + \Delta \dot{x}) \\
&\quad - M\ddot{x}_r - M\Delta \ddot{x} - B\dot{x}_r - B\Delta \dot{x} - Kx_r - K\Delta x\Big]
\end{align}

\textbf{分离参考项和扰动项}：
\begin{align}
\tau &= J^T(\theta)\Big[(\tilde{\Lambda}(\theta)\ddot{x}_r + \tilde{\eta}(\theta, \dot{x}_r) - M\ddot{x}_r - B\dot{x}_r - Kx_r) \\
&\quad + (\tilde{\Lambda}(\theta)\Delta \ddot{x} + \tilde{\eta}(\theta, \dot{x}_r + \Delta \dot{x}) - \tilde{\eta}(\theta, \dot{x}_r) - M\Delta \ddot{x} - B\Delta \dot{x} - K\Delta x)\Big]
\end{align}

\subsection{线性化近似}

\textbf{如果扰动很小}，可以线性化$\tilde{\eta}(\theta, \dot{x}_r + \Delta \dot{x})$：
\begin{equation}
\tilde{\eta}(\theta, \dot{x}_r + \Delta \dot{x}) \approx \tilde{\eta}(\theta, \dot{x}_r) + \frac{\partial \tilde{\eta}}{\partial \dot{x}}\Big|_{\dot{x}_r} \Delta \dot{x}
\end{equation}

\textbf{简化后的控制律}：
\begin{align}
\tau &\approx J^T(\theta)\Big[(\tilde{\Lambda}(\theta)\ddot{x}_r + \tilde{\eta}(\theta, \dot{x}_r) - M\ddot{x}_r - B\dot{x}_r - Kx_r) \\
&\quad + (\tilde{\Lambda}(\theta) - M)\Delta \ddot{x} + \left(\frac{\partial \tilde{\eta}}{\partial \dot{x}}\Big|_{\dot{x}_r} - B\right)\Delta \dot{x} - K\Delta x\Big]
\end{align}

\section{优势和应用}

\subsection{优势}

\begin{enumerate}
\item \textbf{线性化}：
\begin{itemize}
\item 扰动方程是线性的
\item 可以使用线性系统理论
\item 更容易分析和设计
\end{itemize}

\item \textbf{分离关注点}：
\begin{itemize}
\item 参考轨迹跟踪：处理$(\tilde{\Lambda}\ddot{x}_r + \tilde{\eta} - M\ddot{x}_r - B\dot{x}_r - Kx_r)$
\item 扰动抑制：处理$(\tilde{\Lambda}\Delta \ddot{x} + \cdots - M\Delta \ddot{x} - B\Delta \dot{x} - K\Delta x)$
\end{itemize}

\item \textbf{简化计算}：
\begin{itemize}
\item 如果参考轨迹已知，可以预先计算参考项
\item 只需要实时计算扰动项
\item 降低计算复杂度
\end{itemize}
\end{enumerate}

\subsection{应用场景}

\begin{enumerate}
\item \textbf{轨迹跟踪}：
\begin{itemize}
\item 参考轨迹：期望的轨迹
\item 扰动：跟踪误差
\item 目标：使扰动最小化
\end{itemize}

\item \textbf{扰动抑制}：
\begin{itemize}
\item 参考轨迹：标称轨迹
\item 扰动：外力干扰
\item 目标：抑制扰动的影响
\end{itemize}

\item \textbf{稳定性分析}：
\begin{itemize}
\item 分析扰动方程的稳定性
\item 如果扰动方程稳定，系统会回到参考轨迹
\end{itemize}
\end{enumerate}

\section{注意事项}

\subsection{小扰动假设}

\textbf{关键假设}：扰动$\Delta x$, $\Delta \dot{x}$, $\Delta \ddot{x}$必须很小。

\textbf{如果扰动很大}：
\begin{itemize}
\item 线性化近似不成立
\item 需要保留高阶项
\item 或者使用非线性分析
\end{itemize}

\textbf{如何确保扰动小？}
\begin{itemize}
\item 设计好的控制器
\item 使用前馈补偿
\item 快速响应扰动
\end{itemize}

\subsection{参考轨迹的选择}

\textbf{要求}：
\begin{itemize}
\item 参考轨迹应该是可行的（物理上可实现）
\item 参考轨迹应该满足动力学约束
\item 参考轨迹应该平滑（避免大的加速度）
\end{itemize}

\section{总结}

\subsection{替换过程总结}

\begin{enumerate}
\item \textbf{定义参考轨迹}：$x_r(t)$, $\dot{x}_r(t)$, $\ddot{x}_r(t)$

\item \textbf{定义小位移}：
\begin{align}
\Delta x &= x - x_r \\
\Delta \dot{x} &= \dot{x} - \dot{x}_r \\
\Delta \ddot{x} &= \ddot{x} - \ddot{x}_r
\end{align}

\item \textbf{替换}：
\begin{align}
x &= x_r + \Delta x \\
\dot{x} &= \dot{x}_r + \Delta \dot{x} \\
\ddot{x} &= \ddot{x}_r + \Delta \ddot{x}
\end{align}

\item \textbf{代入方程}：将替换代入原始方程

\item \textbf{分离项}：将参考项和扰动项分开

\item \textbf{得到扰动方程}：$M\Delta \ddot{x} + B\Delta \dot{x} + K\Delta x = \Delta f_{\text{ext}}$
\end{enumerate}

\subsection{关键公式}

\begin{enumerate}
\item \textbf{小位移定义}：
\begin{align}
\Delta x &= x - x_r \\
\Delta \dot{x} &= \dot{x} - \dot{x}_r \\
\Delta \ddot{x} &= \ddot{x} - \ddot{x}_r
\end{align}

\item \textbf{扰动方程}：
\begin{equation}
M\Delta \ddot{x} + B\Delta \dot{x} + K\Delta x = \Delta f_{\text{ext}}
\end{equation}

\item \textbf{替换关系}：
\begin{align}
x &= x_r + \Delta x \\
\dot{x} &= \dot{x}_r + \Delta \dot{x} \\
\ddot{x} &= \ddot{x}_r + \Delta \ddot{x}
\end{align}
\end{enumerate}

\subsection{物理意义}

\begin{itemize}
\item \textbf{参考轨迹}：机器人应该遵循的理想轨迹
\item \textbf{小位移}：实际状态与参考状态的偏差
\item \textbf{扰动方程}：描述偏差如何演化
\item \textbf{目标}：使扰动最小化，系统回到参考轨迹
\end{itemize}

\end{document}

